\section{Discusión}
\label{sec:discusion}

El diseño adoptado separa la responsabilidad de cada estructura: la cola de
prioridad por niveles decide el orden de atención, la lista enlazada de
cada nivel permite retirar solicitudes canceladas, y la pila registra el
historial. Esta separación evita resolver todas las necesidades con una
sola estructura, por ejemplo una lista que se recorra linealmente para todo,
lo que degradaría también la atención a $O(n)$.

La principal limitación conocida es el costo lineal de la búsqueda por
identificador, que se traslada a la cancelación y a la verificación de
identificadores repetidos al registrar. Para el tamaño de datos esperado
(una fila de atención, no miles de solicitudes simultáneas) este costo se
considera aceptable en esta etapa. Una mejora prevista es mantener un
índice por identificador que permita ubicar cada solicitud en tiempo
esperado constante sin alterar el orden de atención.

Otra limitación es que el costo constante de la cola de prioridad depende
de que la urgencia tome solo cinco valores. Si en el futuro la prioridad
dependiera de una magnitud continua, como el tiempo de espera acumulado,
habría que evaluar otra representación de la cola de prioridad.

Finalmente, los datos existen solo durante la ejecución del programa, lo
cual es consistente con el alcance definido
(Sección~\ref{sec:metodologia}).
